% Options for packages loaded elsewhere
\PassOptionsToPackage{unicode}{hyperref}
\PassOptionsToPackage{hyphens}{url}
\documentclass[
]{article}
\usepackage{xcolor}
\usepackage{amsmath,amssymb}
\setcounter{secnumdepth}{-\maxdimen} % remove section numbering
\usepackage{iftex}
\ifPDFTeX
  \usepackage[T1]{fontenc}
  \usepackage[utf8]{inputenc}
  \usepackage{textcomp} % provide euro and other symbols
\else % if luatex or xetex
  \usepackage{unicode-math} % this also loads fontspec
  \defaultfontfeatures{Scale=MatchLowercase}
  \defaultfontfeatures[\rmfamily]{Ligatures=TeX,Scale=1}
\fi
\usepackage{lmodern}
\ifPDFTeX\else
  % xetex/luatex font selection
\fi
% Use upquote if available, for straight quotes in verbatim environments
\IfFileExists{upquote.sty}{\usepackage{upquote}}{}
\IfFileExists{microtype.sty}{% use microtype if available
  \usepackage[]{microtype}
  \UseMicrotypeSet[protrusion]{basicmath} % disable protrusion for tt fonts
}{}
\makeatletter
\@ifundefined{KOMAClassName}{% if non-KOMA class
  \IfFileExists{parskip.sty}{%
    \usepackage{parskip}
  }{% else
    \setlength{\parindent}{0pt}
    \setlength{\parskip}{6pt plus 2pt minus 1pt}}
}{% if KOMA class
  \KOMAoptions{parskip=half}}
\makeatother
\usepackage{longtable,booktabs,array}
\newcounter{none} % for unnumbered tables
\usepackage{calc} % for calculating minipage widths
% Correct order of tables after \paragraph or \subparagraph
\usepackage{etoolbox}
\makeatletter
\patchcmd\longtable{\par}{\if@noskipsec\mbox{}\fi\par}{}{}
\makeatother
% Allow footnotes in longtable head/foot
\IfFileExists{footnotehyper.sty}{\usepackage{footnotehyper}}{\usepackage{footnote}}
\makesavenoteenv{longtable}
\setlength{\emergencystretch}{3em} % prevent overfull lines
\providecommand{\tightlist}{%
  \setlength{\itemsep}{0pt}\setlength{\parskip}{0pt}}
\usepackage{bookmark}
\IfFileExists{xurl.sty}{\usepackage{xurl}}{} % add URL line breaks if available
\urlstyle{same}
\hypersetup{
  hidelinks,
  pdfcreator={LaTeX via pandoc}}

\author{}
\date{}

\begin{document}

\section{MetaWingman: a question-first, step-verified, self-improving
agent for systematic review and
meta-analysis}\label{metawingman-a-question-first-step-verified-self-improving-agent-for-systematic-review-and-meta-analysis}

\textbf{Manuscript type:} Methods article (draft for journal-style
adaptation) \textbf{Target style:} Frontiers in Immunology, Methods
section \textbf{Citation style:} Vancouver (numbered citations, order of
first appearance) \textbf{Language:} English

\begin{center}\rule{0.5\linewidth}{0.5pt}\end{center}

\subsection{Abstract}\label{abstract}

\textbf{Background.} Systematic reviews and meta-analyses (SR/MAs)
remain labor-intensive and error-prone despite mature reporting
standards. Recent AI research agents automate pieces of the pipeline,
but few bind the review to a derived question, verify work at the level
of individual methodological steps, and leave an auditable record
through which the system improves itself.

\textbf{Methods.} We built MetaWingman, a five-component agent: a Review
Question Certificate that derives a falsifiable question through a
seven-stage process with hard and soft gates; a ten-stage Socratic
checklist that gates stage entry; a step-level verifier for appraisal
dossiers with abstention and a human-review window; a JSONL audit log
and meta-update loop through which source-carrying proposals are applied
after review; and retrieval components for section-role classification
and evidence retrieval. Components were trained on 109,028
deterministically weak-labeled examples. Validation proceeds through a
six-level ladder (VAL-1, 2a, 2b1, 2b2, 2c, 3) ending in time-sliced
reconstruction of published reviews.

\textbf{Results.} Section-role classification reached macro-F1 0.9995
and evidence retrieval MRR 0.962 (P@1 0.933). Open-set retrieval favored
BM25 single-stage (MRR 0.2649) over trained rerankers, which hurt
two-stage performance. A deterministic R pipeline reproduced a published
random-effects meta-analysis (pooled MD 2.865 vs 2.9; I² 92.67\% vs
93\%) in three locked runs. A first AI-only screening pilot on 649
frozen records achieved gold recall 0.765 (114/149), with abstention the
dominant loss. Cross-provider agreement reached kappa 0.8472.

\textbf{Conclusions.} MetaWingman demonstrates a question-first,
step-verified architecture with reproducible reconstruction evidence.
Current component metrics support consistency with weak labels, not
scientific validity; an AI-only end-to-end review remains unfinished.

\textbf{Keywords:} systematic review; meta-analysis; AI agent; question
certificate; step-level verification; living systematic review

\subsection{中文摘要}\label{ux4e2dux6587ux6458ux8981}

\textbf{背景。} 系统综述与 meta
分析虽有成熟的报告标准，仍耗时且易错。近年的科研 AI
智能体自动化了流程的若干环节，但很少有系统把综述绑定到一个经过推导的问题上、在方法学步骤层面做验证，并留下可审计记录使系统自我改进。

\textbf{方法。} 我们构建了
MetaWingman，一个五组件智能体：七阶段推导且带硬/软门的综述问题证书；十个阶段的苏格拉底清单门禁；带弃权与人工审核窗口的评价步骤级验证器；带出处的
JSONL 审计日志与元更新回路；以及 section-role 分类与证据检索组件。组件在
109,028 条确定性弱标签样本上训练；验证经由六级阶梯（VAL-1 至
VAL-3），终点是已发表综述的时间切分重建。

\textbf{结果。} section-role 分类 macro-F1 0.9995，证据检索候选集 MRR
0.962（P@1 0.933）。开集检索以 BM25 单阶段（MRR
0.2649）为优，训练重排器在两段式配置下为负贡献。确定性 R
管线以三次锁定运行复现了已发表的随机效应 meta 分析（合并 MD 2.865 对
2.9；I² 92.67\% 对 93\%）。首个 AI-only 筛选 pilot（649
条冻结记录）黄金召回 0.765，弃权为主要损失；跨提供方一致性 kappa
0.8472。

\textbf{结论。} MetaWingman
展示了提问先行、步骤验证、可审计自改进的架构，并附复现实证。当前组件指标支持的是与弱标签/准则标签的一致性，而非科学有效性；AI-only
的端到端综述仍未完成。

\textbf{关键词：} 系统综述；meta 分析；AI
智能体；问题证书；步骤级验证；living systematic review

\begin{center}\rule{0.5\linewidth}{0.5pt}\end{center}

\subsection{1. Introduction}\label{introduction}

Systematic review and meta-analysis is the backbone of evidence-based
medicine, yet it is slow, expensive, and sensitive to process errors at
every stage, from search construction to risk-of-bias judgment {[}1{]}.
The conduct standards are mature: the Cochrane Handbook for Systematic
Reviews of Interventions defines the methods {[}1{]}, PRISMA 2020
governs reporting {[}2{]}, PRISMA-S extends reporting to the search
{[}3{]}, and living systematic review frameworks specify how a review
should be updated rather than frozen {[}4{]}. What remains unresolved is
execution. Each of the ten stages a review passes through, topic,
protocol, search, screening, extraction, appraisal, analysis, writing,
reproducibility, and update, is a place where undocumented shortcuts
enter, and where an AI assistant can either compound or contain the
damage.

The practical burden is not just volume. A review is a chain of
judgments in which every downstream decision inherits the assumptions of
the ones before it. A search that misses a database, a screening rule
that silently excludes a population, an extraction that transcribes the
wrong arm, and an appraisal that over-weights a domain all move the
final estimate, and none of them is visible in the published forest
plot. The standards documents tell a review team what must be done; they
do not detect when a step was skipped. This asymmetry, between
well-specified methods and weakly monitored execution, is the specific
problem an agent for systematic review must solve, and it shapes the
design choices in this paper.

The past two years have produced a wave of research agents. The AI
Scientist drafts and reviews machine-learning papers end to end {[}5{]},
AI Scientist-v2 adds a multi-module architecture with self-correction
{[}6{]}, and Google's AI co-scientist frames discovery as a
generate-debate-evolve loop over hypotheses {[}7{]}. Agent Laboratory
targets research workflows including systematic review {[}8{]},
PRefLexOR introduces recursive preference optimization for scientific
reasoning {[}9{]}, and Reflexion {[}10{]}, process verifiers {[}11{]},
and later reflection architectures {[}12{]} established the general
pattern of an agent that inspects its own outputs. MetaSyn explicitly
targets systematic review and meta-analysis as an agentic pipeline
{[}13{]}, and FirstResearch contributes a question-certificate mechanism
that derives and gates review questions before work begins {[}14{]}.

What these systems share is a gap at the level of the review's
scientific contract. First, the question is usually treated as input,
not as a derived, gated artifact; a review built on an underexamined
question cannot be rescued by better execution. Second, verification,
where present, operates on finished outputs (a manuscript, a hypothesis
list), not on individual methodological steps such as an appraisal
judgment; a wrong intermediate step can propagate silently into a
correct-looking synthesis. Third, self-improvement claims are rarely
backed by an auditable trail that ties a proposed change to the failure
that motivated it. Fourth, and most consequential for the methods
literature, capability claims rest on demos rather than on
reconstruction: re-running a published review from its own data and
comparing against the published numbers, which is the closest thing to a
ground truth that the field has.

We report MetaWingman, an agent built on five components that address
these gaps directly: a Review Question Certificate migrated from
FirstResearch {[}14{]}; a ten-stage Socratic checklist that gates stage
entry; a step-level verifier for appraisal dossiers; a JSONL audit log
with a meta-update loop; and retrieval components for section-role
classification and evidence retrieval. We describe the training pipeline
built on deterministically weak-labeled data, a six-level validation
ladder, and time-sliced reconstruction of published reviews. We report
measured results from execution receipts: component metrics, an open-set
retrieval comparison that argues against trained rerankers, a locked
three-run reproduction of a published meta-analysis, a first AI-only
screening pilot, and cross-provider agreement. We do not claim validated
equivalence with human review; the evidence we present supports
consistency with rubric-grounded labels and deterministic reproduction
within a defined engine boundary.

\subsection{2. Methods}\label{methods}

\subsubsection{2.1 System overview}\label{system-overview}

MetaWingman is organized as five components whose outputs are written to
a project control plane (review state, protocol, event ledger) and
validated before they may alter scientific artifacts. The five
components are (1) the Review Question Certificate (RQC), (2) the
ten-stage Socratic checklist, (3) the step-level verifier, (4) the audit
log and meta-update loop, and (5) the retrieval components. Table 1
summarizes each component, the mechanism it implements, and the evidence
reported in this paper.

The control plane separates three kinds of state that are easy to
conflate in an agentic pipeline: the review's protocol and derived
artifacts (what the review claims to do), the event ledger (what the
system actually did, with input and output hashes), and the human
responsibility record (which decisions a person signed). A component's
output enters the review only after the relevant validator accepts it; a
failed validator produces a blocked or abstained event, never silent
continuation. This separation is what makes the reconstruction
experiments in Section 3 possible: because every run is logged with
hashes and frozen configuration, a reproduction claim can be checked
against the log rather than trusted from a narrative.

\subsubsection{2.2 Review Question
Certificate}\label{review-question-certificate}

The RQC is migrated from FirstResearch {[}14{]} and implements the
question-first principle: the review question is derived, not assumed.
The generator proceeds through seven stages that move from primitives to
a falsifiable hypothesis: (i) extraction of primitives (population,
intervention/exposure, comparator, outcome, setting), (ii) statement of
first-principle assumptions, (iii) construction of a mechanism model
linking the question to the evidence landscape, (iv) articulation of the
underlying tension the review will resolve, (v) derivation of a
falsifiable hypothesis, (vi) specification of a minimal decisive test,
and (vii) a novelty gate. The novelty gate is a hard gate backed by a
Europe PMC novelty search plus a soft gate for answerability and
decision relevance; the certificate schema carries both gate types and a
failure-update rule stating what evidence would force revision of the
question. Generation uses max\_tokens 8192 (evidence:
\texttt{metawingman/scripts/generate\_review\_question\_certificate.py}).
To assess the certificate independently of the generator's self-report,
we ran a double-judge blind evaluation in the LLM-as-a-judge paradigm
{[}15{]}: two judge models scored certificates with the audit,
quality-scores, and gate fields stripped, agreement was measured with
Pearson correlation, and the judges were DeepSeek 4.0 and GLM 4.4 in a
smoke-test configuration (evidence:
\texttt{docs/architecture/final-status-2026-08-18.md},
\texttt{docs/architecture/innovation-whitepaper-2026-08-18.md}).

\subsubsection{2.3 Ten-stage Socratic
checklist}\label{ten-stage-socratic-checklist}

Ten checklists, one per stage (topic, protocol, search, screening,
extraction, appraisal, analysis, writing, reproducibility, update), each
contain ten questions (nine mandatory, one optional). A stage may not
begin until the mandatory questions are answered and the checklist gate
passes; the gate is enforced by a deterministic script,
check\_socratic\_checklist.py, which has passed full regression. The
checklist content draws on the Cochrane Handbook {[}1{]}, PRISMA 2020
{[}2{]}, PRISMA-S {[}3{]}, the MECIR conduct standards, the GRADE
Handbook, the living systematic review methods of Elliott et
al.~{[}4{]}, and the agreement interpretation bands of Landis and Koch
{[}16{]} where kappa thresholds are referenced. The purpose of the gate
is procedural, not epistemic: it forces the human or agent to state,
before entering a stage, what the stage will decide and what evidence
would change the decision.

The stage-specific questions are deliberately concrete rather than
generic. The topic checklist asks about the nearest existing review and
the coverage gap relative to it, about the time-window volume of
evidence, and about the falsifiable contribution sentence; the search
checklist asks for source-specific strategies, exact dates, exports, and
query text; the appraisal checklist asks whether every synthesized
result has the required appraisal and whether abstentions were resolved.
Because the checklists are JSON files with a fixed schema, they are
versioned like code, and the gate script can be regression-tested, which
is what allows the project to claim that all ten stages pass the same
deterministic gate rather than a sample of them.

\subsubsection{2.4 Step-level verifier}\label{step-level-verifier}

Verification in MetaWingman operates at the level of individual
methodological steps, not only finished artifacts. The appraisal stage
is guarded by verify\_appraisal\_steps.py, a rule-based verifier
covering ten appraisal steps; a step that fails a required check, or
that lacks a pending human signature, triggers abstention or opens the
human review window, never silent continuation. Appraisal judgments
themselves are produced by a six-domain risk-of-bias classification
component (domains: selection, detection, attrition, reporting,
performance, other), which is described under training below.

The verifier is intentionally rule-based. A learned verifier would
itself need verification, and the project's evidence standard requires
that a failed gate be explainable by a deterministic rule, not by a
model's confidence. The ten steps mirror the appraisal dossier
structure: each dossier must carry the appraisal tool used, the judgment
per domain, the supporting anchor (a quote or table location in the
report), the result-level scope of the judgment, and the reviewer
signature or its explicit absence. The abstention path matters as much
as the pass path: a step that cannot be verified is recorded as
abstained, and the human review window is the only route that turns an
abstention into a decision. This is the mechanism that produced the
abstention-heavy profile observed in the VAL-3 pilot (Section 3.6).

\subsubsection{2.5 Audit log and meta-update
loop}\label{audit-log-and-meta-update-loop}

Every deviation, failure, fix, reflection, question-answer pair, and
certificate update is recorded as a JSONL event (event types: deviation,
failure, fix, reflection, question\_answer, certificate\_update) in the
audit log. Self-improvement is routed through a meta-update loop: a
proposal for change must carry its source, is applied only after passing
a human review window, and the applied entry records the commit that
implemented it. The loop is the mechanism by which the agent's own
failures become training signal for its configuration; at the time of
writing this paper, eight audit-log entries had been recorded (Section
3.8).

Two design rules keep the loop from becoming a rubber stamp. First, a
proposal without a source is not a proposal: the audit-log schema
requires an evidence\_source field that names the failing artifact or
measurement, and entries without one are rejected. Second, the loop
applies to the system's configuration files and documents, not to the
review's scientific state; a proposal can change how the next review is
conducted, but it cannot silently rewrite an already-frozen protocol. In
practice the recorded closures include model-registry corrections after
a live API smoke test, a training-weights fix discovered from a crash,
and the open-retrieval architecture verdict that replaced a provisional
design note. These are small, honest changes, and the paper reports them
as such rather than as evidence of general self-improvement.

\subsubsection{2.6 Retrieval components}\label{retrieval-components}

Two retrieval components serve the search and screening stages. The
section-role classifier assigns a section role to each passage of a
retrieved document so that screening and extraction can target the parts
of a report that carry the relevant evidence. The evidence-retrieval
component selects candidate passages for extraction. Both are trained
models; their measured performance is reported in Section 3.2, and an
open-set retrieval comparison that motivated keeping BM25 as the
open-retrieval baseline is reported in Section 3.3.

The two components answer different questions and are evaluated
differently. Section-role classification is a closed-set labeling task:
each passage receives one role, and macro-F1 against the labeled
evaluation split is the natural metric. Evidence retrieval is a ranking
task within a candidate set: given a review's field-plus-title query,
rank the supporting passages, measured by MRR and P@1 on frozen
hard-negative pairs. The distinction between candidate-set ranking and
open retrieval is not cosmetic; Section 3.3 shows that the trained
retrieval model, which ranks well inside a curated candidate set,
performs at noise level when asked to retrieve from the full corpus. The
design therefore treats the trained components as refiners of a lexical
baseline, never as the first-stage retriever.

\subsubsection{2.7 Training pipeline}\label{training-pipeline}

Training data consists of 109,028 examples across two tasks (54,514 per
task), split into train (87,264) and dev (21,764); all labels are
deterministic weak labels produced by rules, not by models, which makes
the labels reproducible by construction (evidence:
\texttt{docs/architecture/training-run-report-2026-08-17.md}). The
appraisal six-domain component has three generations. The V3 model is
trained on rule weak labels and scored on held-out appraisal passages
(dev macro-F1 0.8500), a measurement of consistency with the rule
labels, not of correctness against a reference standard. A
criterion-level evaluation (VAL-2c) scored 100 appraisal criteria
against the task's guideline and found only slight agreement (kappa
0.311, 95\% CI 0.191 to 0.431), which we interpreted as a mismatch
between the judgment rules and the guideline's criteria; the component
was then re-labeled under criterion supervision, yielding 9,906
re-labeled records (other 8,443, selection 617, reporting 488,
performance 145, attrition 112, detection 101). A rubric-based
generation (rubric V2) reached dev macro-F1 0.3777 and weighted-F1
0.871, and a generation-2 spot check on 100 items reached agreement 0.80
with weighted-F1 0.871, with kappa not statistically meaningful because
98\% of items fell in the majority class (other). Wherever we report
these numbers, the claim is consistency with the corresponding labels or
rubric, not external validity.

The decision to use deterministic weak labels rather than
model-generated labels deserves emphasis. Model-generated labels carry
the annotator's own biases into the training signal, and a model trained
on them reproduces those biases with high agreement; the project's data
policy therefore treats any provider-validated JSON object as a
candidate label, never as gold. Deterministic rules have the opposite
property: they are cheap to audit, versionable, and their failure modes
are legible, which is exactly what the VAL-2c evaluation exploited when
it found the rules at odds with the guideline. The cost is that
rule-consistency ceilings bound what the trained models can reach, and
the rubric V2 macro-F1 of 0.3777 is the visible consequence of
relabeling against a harder reference: the model now tracks the rubric's
distinctions, but the majority class still dominates the macro average.

\subsubsection{2.8 Validation ladder}\label{validation-ladder}

Validation is organized as a six-level ladder, each level answering a
different question:

\begin{itemize}
\tightlist
\item
  VAL-1: select and license-review published-review reconstruction
  families without contaminating held-out data (first promotion: an
  analysis-slice case from Hodgkiss et al.~{[}18{]}, BSD-3-Clause
  materials; remaining families guard-test-only or blocked, recorded per
  family);
\item
  VAL-2a: encode the AI-only repeated-run plan, the
  published-expert/corrected-reference policy, integrity exclusions, the
  error taxonomy, the metric set, and inference limits;
\item
  VAL-2b1: freeze the component-axis task manual --- scientific loss
  weights, release thresholds, configuration and prompt hashes, and
  stopping rules --- before any held-out model run;
\item
  VAL-2b2: fill and freeze reconstruction-case task manuals (first case
  manual sealed for the analysis slice; remaining families open);
\item
  VAL-2c: freeze a 100-item appraisal-domain spot-check with a sealed
  answer key and a SHA-256 manifest, and score the weak-label rules
  against rubric-grounded judgment (kappa 0.311, Section 3.1);
\item
  VAL-3: run the AI-only configuration/ablation pilot (first AI-only
  screening run reported in Section 3.6; full configuration ladder still
  open).
\end{itemize}

Agreement metrics use Cohen's kappa {[}17{]}, interpreted with the
Landis and Koch bands {[}16{]} introduced in Section 2.3; model-vs-model
and judge-vs-judge comparisons are reported as Pearson correlations or
kappa as stated per experiment. The ladder is deliberately ordered so
that the weakest form of evidence (execution) precedes the strongest
form currently available (deterministic reconstruction of published
numbers); no level above VAL-1 is claimed until the lower levels pass.

Two properties of the ladder matter for how its results should be read.
First, each level measures agreement with a different reference, and the
references are not interchangeable: agreement with rule labels does not
imply agreement with the guideline, which is precisely the failure the
VAL-2c result exposed. Second, the ladder does not contain an
independent human-reference level for the appraisal component; the
criterion-supervised relabeling of VAL-2b2 is a re-anchoring to the
guideline, not an independent human judgment. This is a stated
limitation of the current evidence, not an oversight, and it is repeated
in the Discussion.

\subsubsection{2.9 Reconstruction
mechanism}\label{reconstruction-mechanism}

Reconstruction follows the living systematic review principle of time
slicing {[}4{]}: the published review's search cutoff is honored, and
only records available before that cutoff are used. For the first
reconstruction case, the deterministic analysis pipeline is a locked R
pipeline (metafor engine) run three times under identical configuration;
the three runs must agree byte-for-byte on outputs before scoring
against the published numbers. Scoring uses declared tolerances (pooled
effect ±0.05, CI bounds ±0.05, I² ±1.0 percentage point, study count
exact, Egger significance side). For the ag-rdt living-update family,
the search strategies were extracted verbatim from the published
supplementary materials, and the later review's workbook was used as a
read-only reference; the operational corpus is a merged, deduplicated
slice of PubMed records and preprints assembled under a documented
protocol (Section 3.5).

Reconstruction is the project's answer to the demo problem identified in
the Introduction. A capability claim about meta-analysis is hard to
falsify on a synthetic example, because the correct answer is whatever
the pipeline produced; a published review with published numbers
supplies an external anchor. The time-slice rule keeps the exercise
honest in the other direction: the reconstruction cannot cheat by using
records that were not available to the original team. The same logic
extends to the living-update family, where the reference is not a single
number but a sequence of included-study sets across versions, and where
the scoring tolerances are frozen before the reference is unsealed.

\subsubsection{2.10 Ethics and reporting}\label{ethics-and-reporting}

This study uses only published open-access literature and previously
reported data; no human subjects were involved. All experiments were
preregistered or protocol-frozen where the ladder requires it (VAL-3 in
particular). Reporting follows the spirit of PRISMA 2020 {[}2{]} and
PRISMA-S {[}3{]} for the retrieval components, adapted for a methods
paper.

\subsection{3. Results}\label{results}

All numbers in Section 3 come from machine execution receipts or the
project's result documents; each number's first occurrence names its
evidence source. The receipts and result documents are listed in Data
Availability.

\subsubsection{3.1 Training corpus and component
generations}\label{training-corpus-and-component-generations}

The frozen training corpus contains 109,028 weakly supervised examples
(87,264 train, 21,764 dev), all deterministically labeled (evidence:
\texttt{docs/architecture/training-run-report-2026-08-17.md}). Table 2
reports the appraisal six-domain component across its three generations.
The V3 rule-weak-label model reached dev macro-F1 0.8500 (evidence:
\texttt{docs/architecture/appraisal-step-component-results-2026-08-18.md});
the rubric V2 generation reached macro-F1 0.3777 and weighted-F1 0.871
(evidence:
\texttt{docs/architecture/appraisal-rubric-v2-results-2026-08-18.md});
the generation-2 spot check on 100 items achieved agreement 0.80 and
weighted-F1 0.871, with kappa uninformative under the
98\%-majority-class distribution (evidence:
\texttt{docs/architecture/appraisal-step-component-results-2026-08-18.md}).
The criterion-level evaluation (VAL-2c) scored 100 criteria against the
guideline and produced kappa 0.311 (95\% CI 0.191 to 0.431), interpreted
as a rule-vs-guideline mismatch (evidence:
\texttt{docs/architecture/val2c-scoring-result-2026-08-18.md}); the
resulting criterion-supervised relabeling produced 9,906 records whose
class distribution (other 8,443, selection 617, reporting 488,
performance 145, attrition 112, detection 101) is heavily skewed toward
the majority class, which caps the informativeness of macro-averaged
metrics (evidence:
\texttt{docs/architecture/appraisal-rubric-v2-results-2026-08-18.md};
decision in
\texttt{docs/architecture/appraisal-task-relabeling-decision-2026-08-18.md}).

Two readings of Table 2 are worth making explicit. The V3 macro-F1 of
0.8500 is high because the model learns the rules that produced its
labels, and the rule-labels are internally consistent by construction;
that number says nothing about whether the rules match the guideline,
which is the question VAL-2c answered with kappa 0.311. After
relabeling, the rubric V2 model tracks the harder reference, and its
weighted-F1 of 0.871 with a macro-F1 of 0.3777 reflects a task in which
85.2\% of records fall in the other class (8,443 of 9,906); a model can
be right on the bulk of the data while remaining weak on the rare,
decision-relevant domains. The generation-2 spot check (agreement 0.80)
is reported with its caveat: with 98\% of items in the majority class,
agreement is dominated by the majority, and kappa is not a meaningful
summary.

\subsubsection{3.2 Retrieval components}\label{retrieval-components-1}

The section-role classifier reached eval macro-F1 0.9995 on its
evaluation split (evidence:
\texttt{docs/architecture/training-run-report-2026-08-17.md}). The
evidence-retrieval component reached candidate-set MRR 0.962 and P@1
0.933 (evidence:
\texttt{docs/architecture/training-run-report-2026-08-17.md};
receipt-verified in
\texttt{validation-output/server-download/evidence-retrieval-execution-receipt.json}
for loss and recall). These numbers are candidate-set metrics: they
describe ranking within a provided candidate set, not open retrieval.

The section-role result at 0.9995 macro-F1 should be read as
near-ceiling performance on a closed labeling task with a well-specified
label scheme; it does not mean that document structure is understood in
any general sense, only that the role of a passage is almost always
recoverable from its text and position. The evidence-retrieval numbers
likewise describe a narrow job, ranking the correct supporting passage
among a small set of hard negatives built for the review's
field-plus-title query. Both metrics are reported because they are the
ones the components were designed and trained for; the open-set question
is handled separately in Section 3.3.

\subsubsection{3.3 Open-set retrieval
decision}\label{open-set-retrieval-decision}

Because the candidate-set numbers do not transfer to open retrieval, we
benchmarked single-stage and two-stage retrieval on the full operational
corpus (evidence:
\texttt{docs/architecture/bm25-two-stage-results-2026-08-18.md} and
\texttt{docs/architecture/retrieval-direction-evidence-2026-08-18.md}).
BM25 single-stage (k1=1.5, b=0.75) reached MRR 0.2649, above TF-IDF
(0.2199) and far above the trained 110M-parameter model (0.0045), which
performed at noise level on open queries. The trained reranker was
negative in the two-stage configuration: two-stage MRR fell
monotonically with candidate-set size K (0.128 at K=50 to 0.058 at
K=200). BM25 recall ceilings were 0.4300, 0.4655, and 0.5071 at K=50,
100, and 200. Two conclusions follow. Open retrieval in MetaWingman uses
BM25 single-stage as the baseline, and the trained reranker is reserved
for candidate-set refinement where its MRR 0.962 applies. The 0.507
recall ceiling at K=200 also quantifies how much evidence is unreachable
by lexical retrieval alone, a bound the search stage must disclose.

The comparison is deliberately blunt. The trained model's open-retrieval
MRR of 0.0045 is below random ranking for practical purposes, which
rules out using it as a first-stage retriever; and the two-stage decline
with K shows that the reranker's errors accumulate as more distractors
enter the candidate set, exactly the regime that open retrieval
produces. The recall ceiling deserves the same bluntness: if lexical
retrieval cannot surface half of the relevant records even at K=200,
then a search built on lexical anchors alone is structurally incomplete,
and the screening stage must either widen the anchor vocabulary or
accept the bound with a documented rationale. These are design
constraints measured on one corpus, and the project treats them as
corpus-level evidence rather than universal laws.

\subsubsection{3.4 First reconstruction case: Hodgkiss et
al.~2023}\label{first-reconstruction-case-hodgkiss-et-al.-2023}

The first time-sliced reconstruction case is the random-effects
meta-analysis of peak oxygen uptake (V̇O2peak) from Hodgkiss et
al.~{[}18{]} (doi:10.1371/journal.pmed.1004082), using the review's
BSD-3-Clause materials and its time slice. The deterministic R pipeline
(metafor engine) ran three times under locked configuration; all three
runs produced identical outputs, and scoring passed every declared
tolerance (evidence:
\texttt{validation-output/reconstruction-runs/sci-exercise/rep-1/execution-receipt.json}
and
\texttt{docs/architecture/reconstruction-first-case-results-2026-08-18.md}).
Table 3 reports the comparison. The pooled mean difference was 2.865
versus the published 2.9 (tolerance ±0.05), the CI was 1.795 to 3.935
versus 1.8 to 3.9, I² was 92.67\% versus 93\% (±1.0 percentage point),
the study count matched exactly (k=16), and the Egger test was
non-significant on the same side as published (0.483 versus 0.54). The
declared boundary matters: this is a recomputation within the same
metafor engine family, which demonstrates that our pipeline reproduces
published analysis mechanics, not that any model performed the analysis.

Why does this case count as evidence at all if it only re-runs a
published analysis? Because the pipeline that produced it is the same
pipeline the project uses for the analysis stage, with the same input
staging, hash verification, locked environment, and scoring code. A
reproduction that passes all five scored dimensions on a real published
review, in three byte-identical locked runs, establishes that the
mechanical core of the analysis stage is sound: the data staging does
not corrupt values, the effect-size computation matches the published
engine's conventions, and the tolerance scheme is calibrated to real
published precision. What it does not establish is any claim about model
capability, since no model was involved; that is why the case is
presented as a pipeline-mechanics result and why the AI-only end-to-end
experiment remains the decisive unfinished test.

\subsubsection{3.5 Living-update family: ag-rdt (Brümmer et
al.)}\label{living-update-family-ag-rdt-bruxfcmmer-et-al.}

The second reconstruction family covers the living updates of a
systematic review of SARS-CoV-2 antigen rapid diagnostic test (Ag-RDT)
accuracy, published by Brümmer et al.~in 2021 {[}19{]} and updated in
2022 {[}20{]}; the 2021 version was subsequently corrected {[}21{]}. The
2021 review's search strategies were extracted verbatim from the
published supplementary (S2 Text) with a date-window correction, and the
2022 workbook (CC BY-NC-ND 4.0) was used as a read-only comparison
reference (evidence:
\texttt{docs/architecture/ag-rdt-living-update-case-design-2026-08-18.md}).
The operational corpus is a merged slice of 12,498 records: 7,500 PubMed
records retrieved via NCBI-native endpoints plus 5,000 preprint records,
deduplicated under a documented protocol (evidence:
\texttt{validation-output/ag-rdt-corpus/merged/execution-receipt.json}).
A deterministic screening-chain rehearsal on this corpus showed that
lexical term anchors retained 42\% of records, which we classify as
triage-level retention, and abstract-level extraction coverage was
limited (sensitivity 17.4\%, specificity 14.7\%, n 41.5\%) (evidence:
\texttt{docs/architecture/ag-rdt-living-update-case-design-2026-08-18.md};
receipts in
\texttt{validation-output/ag-rdt-corpus/screened-rehearsal-v2/execution-receipt.json}).
The low extraction coverage quantifies the ceiling of abstract-only
processing and motivates full-text retrieval for extraction, a step the
pipeline documents as required.

The living-update family tests a different property than the first case:
not whether a single pooled estimate can be reproduced, but whether the
machinery of versioned evidence, search reuse, and corpus assembly works
when the reference moves across versions. The corpus numbers are the
measurable part of that test. The 42\% retention of the lexical anchors
is an order of magnitude above the review's real inclusion rate, which
is exactly the property a triage filter should have: it keeps almost
everything that matters while cutting the pool by more than half. The
abstract-level extraction coverage (sensitivity 17.4\%) is the
counterweight: triage at the record level is not extraction at the field
level, and the deterministic chain is documented as triage-plus-audit,
with field-level extraction deferred to the full-text stage or the
AI-assisted arm. The 2022 workbook was used read-only because its CC
BY-NC-ND 4.0 license forbids redistribution; the comparison reference is
therefore cited, not copied.

\subsubsection{3.6 VAL-3: first AI-only screening
pilot}\label{val-3-first-ai-only-screening-pilot}

The first VAL-3 AI-only screening run used 649 frozen records, of which
149 were gold matches (merged-corpus matches of the 2022 update's
included studies) and 500 were non-gold (evidence:
\texttt{docs/architecture/val3-ai-screening-pilot-results-2026-08-18.md}).
Gold recall was 0.765 (114/149); the dominant loss was abstention (26 of
149 gold records, 17.4\%), not false exclusion. The pipeline retained
19.4\% of records, versus the deterministic anchors' 42\% retention, a
substantially more conservative operating point. The pilot ran with a
resume-friendly runner over 649 calls with one schema-repair retry
before abstention. These numbers support feasibility of AI-first
screening with an abstention-heavy profile; they do not establish
equivalence with dual independent human screening.

The composition of the recall loss is the pilot's most informative
finding. Twenty-six of the 149 gold records were abstained rather than
excluded, which means the loss is recoverable in principle: those
records can be routed to the human review window without any
re-screening. False exclusion, the error that cannot be recovered
without re-running the whole screen, was not the binding constraint. The
cost of this design is visible in the retention rate: 19.4\% is roughly
half the deterministic anchors' 42\%, so the AI-first screen keeps a
larger candidate pool for the human step. Whether that trade is
acceptable depends on the review's screening volume and the price of
missing a record; the pilot's numbers make the trade explicit rather
than hidden in a single accuracy figure. One pilot on 649 records is a
feasibility result, and the preregistered protocol for a full AI-only
configuration remains the next step.

\subsubsection{3.7 Cross-provider and judge
agreement}\label{cross-provider-and-judge-agreement}

A cross-provider comparison ran the same blinded set of 999 paragraphs
through a GLM model (glm-5.2 C3) and a DeepSeek model (C3-R2) (evidence:
\texttt{docs/architecture/glm-cross-provider-results-2026-08-18.md}).
Section-role macro-F1 was 0.8816 for GLM (999 scored, 0 abstained after
three dead-letter recovery rounds; the dominant dead-letter cause was
reasoning-token budget exhaustion, all recovered) versus 0.9385 for
DeepSeek; retrieval selection accuracy was 0.96 versus 0.93; and across
providers, Cohen's kappa on the shared 999 scored passages was 0.8472
(95\% CI 0.8221 to 0.8722), interpreted as substantial to almost-perfect
agreement {[}16{]}. These numbers indicate that the two providers
largely agree on section-role assignment, which supports provider
diversity for verification roles without claiming that either model is
correct where both agree. The certificate double-judge smoke test
(DeepSeek 4.0 and GLM 4.4) ran blind with audit, quality-scores, and
gate fields stripped and agreement measured by Pearson correlation
(evidence:
\texttt{docs/architecture/innovation-whitepaper-2026-08-18.md}).

Provider diversity is a verification strategy, not a truth oracle. If
two independent providers assign the same section role to the same
passage, the assignment is more likely to be right than either model
alone, but two models can share the same systematic error, and the kappa
of 0.8472 does not measure correctness against any reference. The point
of the comparison is narrower: the system can run its verification roles
on two providers and expect them to agree most of the time, which makes
cross-provider double-checking a usable mechanism rather than a source
of constant disagreement. The per-metric differences (0.8816 versus
0.9385 on section-role; 0.96 versus 0.93 on retrieval selection) are
also a reminder that provider choice moves component performance by a
few points, which matters when a component sits at a decision boundary.

\subsubsection{3.8 Audit-log closures and system-level
evidence}\label{audit-log-closures-and-system-level-evidence}

At the time of writing, eight audit-log entries had been recorded, of
which seven carried an applied commit through the review window
(evidence: \texttt{validation-output/audit-log.jsonl}, 8 entries; design
in \texttt{docs/architecture/innovation-whitepaper-2026-08-18.md}). The
ten-stage checklist gate (check\_socratic\_checklist.py) passed full
regression across all ten stages (evidence:
\texttt{docs/architecture/final-status-2026-08-18.md}). Table 1 maps
each of the five components to its mechanism and its evidence.

The eight entries are a process record, not a performance claim. They
include corrections to the model registry after a live API smoke test, a
training-weight fix, the open-retrieval verdict, the VAL-2c pivot, and
the GLM integration lesson, each logged with its evidence source and,
where applied, its commit. The honest reading is that the loop works as
a bookkeeping mechanism: failures are recorded, proposals carry sources,
and applications are traceable. Whether the loop produces durable
capability gains is a question that requires longitudinal data the
project does not yet have, and the paper does not claim it.

\subsubsection{3.9 Comparison with
MetaSyn}\label{comparison-with-metasyn}

A task-mapping comparison against MetaSyn {[}13{]} covered 15
review-related tasks: 6 fully covered, 7 partially covered, 2 gaps
(evidence: \texttt{docs/architecture/final-status-2026-08-18.md}; detail
in \texttt{docs/architecture/metasyn-adapter-design-2026-08-18.md}). We
also audited dataset-card licensing for MetaSyn's released assets
(MIT-annotated labels with upstream third-party terms) against MetaSyn's
422 reviews and 140,585-record corpus (evidence:
\texttt{docs/architecture/top-journal-contribution-story.md} for the
corpus identity;
\texttt{docs/architecture/metasyn-adapter-design-2026-08-18.md} for the
license audit); the audit confirmed the license posture is compatible
with research use but that redistribution of third-party text requires
upstream terms.

The mapping is a scoping exercise, not a head-to-head benchmark: it
compares which tasks each system claims to cover, on the basis of
MetaSyn's documentation and our own component inventory, without running
either system on the other's data. The two gaps identified are
informative because they point at machinery MetaSyn does not advertise:
the gold-corpus evaluation axis and stage attribution of metrics, both
of which this paper treats as load-bearing for validation. The license
audit matters for a different reason: MetaSyn's 422-review and
140,585-record corpus is a natural third-party evaluation axis for
retrieval and screening components, and the audit establishes that using
it for research evaluation is compatible with the licenses, while
redistributing its third-party text is not.

\subsection{4. Discussion}\label{discussion}

MetaWingman's contribution is architectural rather than a single
performance number: a question-first derivation step, verification at
the level of individual methodological steps, an auditable
self-improvement loop, and reconstruction evidence tied to published
reviews. Each of the five components addresses a failure mode that the
existing agent literature leaves open. FirstResearch {[}14{]} supplies
the question-certificate mechanism we migrated; our addition is a
double-judge blind evaluation of the certificate and its hard/soft gate
structure. MetaSyn {[}13{]} is the closest system in scope, an agentic
SR/MA pipeline; the task-mapping comparison (6/15 fully covered, 7/15
partial, 2/15 gaps) locates the boundary between the two systems, and
our gap analysis identifies where MetaSyn's coverage is incomplete
relative to the ten-stage contract. Reflexion {[}10{]} and process
verifiers {[}11{]} establish the self-inspection pattern that our
step-level verifier and audit loop instantiate in the appraisal domain;
our verifier is deterministic and rule-based by design, which trades
flexibility for auditability.

The comparison with the wider agent literature is deliberately
conservative. Systems such as the AI Scientist {[}5{]} and AI
Scientist-v2 {[}6{]} demonstrate what end-to-end automation looks like
when the output is a paper and the review loop is a peer-review proxy;
the AI co-scientist {[}7{]} shows a generate-debate-evolve loop over
hypotheses; Agent Laboratory {[}8{]} and PRefLexOR {[}9{]} extend the
pattern to broader research workflows. MetaWingman does not compete with
these systems on coverage. Its claim is narrower and, we argue, more
falsifiable: that the review's scientific contract can be enforced by
derivation, step-level verification, and reconstruction, and that each
of these mechanisms can be measured. Where the existing systems validate
with demos and self-report, this paper validates with published-number
reconstruction and execution receipts; where they improve by reflection
over finished outputs {[}10{]}{[}12{]}, this paper improves by logged
proposals that must carry a source and survive a review window. The two
styles are complementary, and a production system would plausibly want
both.

Three results deserve emphasis because they run against the default
assumptions of the field. First, the trained 110M reranker was worse
than BM25 on open retrieval (MRR 0.0045 versus 0.2649) and actively
harmful in two-stage configuration (0.128 to 0.058 as K grew). Training
on candidate sets teaches discrimination within a candidate
distribution, and that does not transfer to open queries; the paper's
open-retrieval decision (BM25 baseline, trained reranker only in
candidate sets) is a direct consequence of measured data, not of
preference. Second, the VAL-2c kappa of 0.311 (95\% CI 0.191 to 0.431)
between the rule-based appraisal judgments and the guideline's criteria
is the single most informative negative result in the system: it shows
that weak-label agreement (macro-F1 0.8500 against rule labels) is not
evidence of agreement with the guideline, and it drove the
criterion-supervised relabeling of 9,906 records. Third, the
reconstruction of Hodgkiss et al.~{[}18{]} within tolerance on all five
scored dimensions, with three byte-identical locked runs, demonstrates
that the deterministic analysis chain is reproducible at the level of
published numbers; the declared boundary (same metafor engine family) is
what keeps this claim honest.

The limitations are substantial and we state them explicitly. (1) All
component metrics measure consistency with rule-based or rubric-grounded
labels; none of them measures scientific validity against an independent
human reference standard beyond the VAL-3 pilot's gold set. The rubric
V2 macro-F1 of 0.3777 alongside weighted-F1 0.871, and the
98\%-majority-class distribution of the relabeled data, both limit what
macro-averaged scores can tell us. (2) The criterion-level kappa 0.311
documents a rule-vs-guideline mismatch that was repaired by relabeling
but not independently re-verified against human judges; the kappa itself
has a wide confidence interval (0.191 to 0.431) reflecting
evaluation-process variance on a 100-item sample. (3) The reconstruction
case is a recomputation within the same statistical engine family; it
validates pipeline mechanics, not model capability, and generalizes only
as far as metafor-compatible analyses. (4) The VAL-3 pilot is a single
AI-only screening run on 649 frozen records with an abstention-heavy
profile (gold recall 0.765, retention 19.4\%); equivalence with dual
independent human screening is not established, and the abstention loss
(17.4\% of gold) is the main barrier to recall. (5) The abstract-level
extraction coverage (sensitivity 17.4\%) shows that screening and
extraction cannot rely on abstracts, and the full-text retrieval stage
required by the pipeline was not part of the numbers reported here. (6)
An AI-only end-to-end review, from certificate to published-numbers
reconstruction under one protocol, remains unfinished; what we report is
a validated assembly of components plus one locked reconstruction slice.
(7) Cross-provider agreement (kappa 0.8472) is measured on a single
999-paragraph set with one model per provider; it supports provider
diversity as a verification mechanism but says nothing about absolute
correctness. (8) The evaluation processes themselves are not yet
independently replicated: the VAL-2c kappa, the cross-provider kappa,
and the judge correlations each rest on one evaluation run, and the
confidence intervals reported for the kappas quantify sampling variance,
not replication stability.

Where does the project go next? The immediate target is a complete
AI-only end-to-end run on a preregistered case, closing the loop from
certificate through screening, extraction, appraisal, and synthesis to a
frozen analysis that is scored against the published reference under the
VAL-3 protocol. The open problems are the abstention policy (26 of 149
gold records abstained in VAL-3; recall is capped by abstention, not by
false exclusion), the appraisal criterion gap (the kappa 0.311 class of
failure), and the lexical recall ceiling (0.507 at K=200), which
together determine how much of the evidence surface an AI-first pipeline
can reach. Two further directions follow from the limitations: an
independent human-reference study for the appraisal component, which
would give the ladder the level it currently lacks, and a second
reconstruction case outside the metafor engine family, which would test
whether the pipeline-mechanics result generalizes beyond one statistical
engine. Both are preregistration-shaped experiments, and the project's
protocol discipline applies to them as it does to the runs reported
here.

\subsection{5. Conclusion}\label{conclusion}

MetaWingman binds an SR/MA pipeline to a derived, gated question;
verifies appraisal at the level of individual steps; records every
failure and fix in an auditable log with a source-carrying meta-update
loop; and validates its components through a six-level ladder that ends
in time-sliced reconstruction of published reviews. The measured record
supports three claims: deterministic reproduction of a published
meta-analysis within tolerance across three locked runs; component-level
consistency with rule- and rubric-grounded labels; and a first AI-only
screening pilot with abstention-dominated recall loss. It does not yet
support a claim of validated equivalence with human review, and the
end-to-end AI-only review remains the decisive unfinished experiment.
The design goal of the project is that every number in this paper is
reproducible from its execution receipt, which we treat as the minimal
standard for methods claims in this field.

\begin{center}\rule{0.5\linewidth}{0.5pt}\end{center}

\subsection{Data Availability}\label{data-availability}

All evidence cited in this paper is machine-generated and stored in the
project repository under \texttt{validation-output/} (execution
receipts, JSON) and \texttt{docs/architecture/} (result documents,
Markdown). Key receipts: the frozen training corpus report
(\texttt{docs/architecture/training-run-report-2026-08-17.md}),
including section-role macro-F1 0.9995 and evidence-retrieval
candidate-set MRR 0.962/P@1 0.933, with server receipts in
\texttt{validation-output/server-download/section-role-execution-receipt.json}
and \texttt{evidence-retrieval-execution-receipt.json}; appraisal
component results
(\texttt{docs/architecture/appraisal-step-component-results-2026-08-18.md},
\texttt{appraisal-rubric-v2-results-2026-08-18.md},
\texttt{val2c-scoring-result-2026-08-18.md},
\texttt{appraisal-task-relabeling-decision-2026-08-18.md});
open-retrieval benchmarks
(\texttt{docs/architecture/bm25-two-stage-results-2026-08-18.md},
\texttt{retrieval-direction-evidence-2026-08-18.md}); reconstruction
receipts
(\texttt{validation-output/reconstruction-runs/sci-exercise/rep-\{1,2,3\}/execution-receipt.json}
and
\texttt{docs/architecture/reconstruction-first-case-results-2026-08-18.md});
the ag-rdt corpus and rehearsal receipts
(\texttt{validation-output/ag-rdt-corpus/**/execution-receipt.json},
\texttt{docs/architecture/ag-rdt-living-update-case-design-2026-08-18.md});
the VAL-3 pilot
(\texttt{docs/architecture/val3-ai-screening-pilot-results-2026-08-18.md});
the cross-provider comparison
(\texttt{docs/architecture/glm-cross-provider-results-2026-08-18.md});
the certificate double-judge smoke test and the audit-log design
(\texttt{docs/architecture/innovation-whitepaper-2026-08-18.md},
\texttt{docs/architecture/final-status-2026-08-18.md}); the audit-log
closures (\texttt{validation-output/audit-log.jsonl}); and the MetaSyn
task mapping, corpus identity, and license audit
(\texttt{docs/architecture/final-status-2026-08-18.md},
\texttt{docs/architecture/top-journal-contribution-story.md},
\texttt{docs/architecture/metasyn-adapter-design-2026-08-18.md}).
Analysis code and the deterministic R reconstruction pipeline are
versioned in the repository; the reconstruction inputs use the published
review's BSD-3-Clause materials, and the 2022 ag-rdt workbook is used
read-only under its CC BY-NC-ND 4.0 license. Raw PDF/XML corpora are
kept outside Git per the project's data policy.

\subsection{Ethics Declaration}\label{ethics-declaration}

This study did not involve human subjects, patient data, or animal
experiments. It uses published open-access literature, publicly released
benchmark materials, and previously reported data. All reuse respects
the source licenses (BSD-3-Clause for the Hodgkiss materials; CC
BY-NC-ND 4.0 read-only for the 2022 ag-rdt workbook; MIT annotations
with upstream third-party terms for the MetaSyn assets).

\subsection{Author Contributions
(CRediT)}\label{author-contributions-credit}

{[}Placeholder author list; to be completed by the author team.{]}
Conceptualization: {[}Author A{]}; Methodology: {[}Author A{]},
{[}Author B{]}; Software: {[}Author A{]}, {[}Author B{]}; Validation:
{[}Author A{]}, {[}Author B{]}; Formal analysis: {[}Author B{]}; Data
curation: {[}Author B{]}; Writing (original draft): {[}Author A{]};
Writing (review and editing): {[}Author A{]}, {[}Author B{]}; Project
administration: {[}Author A{]}.

\subsection{Competing Interests}\label{competing-interests}

The authors declare no competing interests.

\subsection{Funding}\label{funding}

This work received no specific grant from any funding agency in the
public, commercial, or not-for-profit sectors.

\subsection{AI-use Statement}\label{ai-use-statement}

This manuscript was drafted with AI assistance and verified against the
project's evidence dossiers and execution receipts. All reported numbers
originate from machine execution receipts or versioned result documents
generated by the MetaWingman pipeline; no number was generated during
manuscript writing. The draft follows a prespecified evidence dossier,
and each figure was checked against its cited receipt. The ten-stage
checklist gate and audit-log discipline of the project were applied to
the manuscript itself.

\begin{center}\rule{0.5\linewidth}{0.5pt}\end{center}

\subsection{Tables}\label{tables}

\textbf{Table 1. The five components of MetaWingman, their mechanisms,
and the evidence reported in this paper.}

{\def\LTcaptype{none} % do not increment counter
\begin{longtable}[]{@{}
  >{\raggedright\arraybackslash}p{(\linewidth - 4\tabcolsep) * \real{0.3333}}
  >{\raggedright\arraybackslash}p{(\linewidth - 4\tabcolsep) * \real{0.3333}}
  >{\raggedright\arraybackslash}p{(\linewidth - 4\tabcolsep) * \real{0.3333}}@{}}
\toprule\noalign{}
\begin{minipage}[b]{\linewidth}\raggedright
Component
\end{minipage} & \begin{minipage}[b]{\linewidth}\raggedright
Mechanism
\end{minipage} & \begin{minipage}[b]{\linewidth}\raggedright
Evidence reported
\end{minipage} \\
\midrule\noalign{}
\endhead
\bottomrule\noalign{}
\endlastfoot
Review Question Certificate & Seven-stage derivation (primitives to
falsifiable hypothesis), hard novelty gate (Europe PMC) and soft
answerability gate, failure-update rule; max\_tokens 8192 & Double-judge
blind evaluation (DeepSeek 4.0, GLM 4.4), Pearson agreement; hard/soft
gate structure \\
Ten-stage Socratic checklist & Ten stages × ten questions (9 mandatory +
1 optional); deterministic gate script & Full regression pass of
check\_socratic\_checklist.py; sources: Cochrane Handbook, PRISMA 2020,
PRISMA-S, MECIR, GRADE Handbook, Elliott et al. \\
Step-level verifier & verify\_appraisal\_steps.py: ten appraisal steps,
abstention and human-review window & Six-domain RoB classification
component (Table 2) \\
Audit log + meta-update loop & JSONL events
(deviation/failure/fix/reflection/question\_answer/certificate\_update);
source-carrying proposals applied via review window with commit recorded
& Eight audit-log entries recorded, seven applied (Section 3.8) \\
Retrieval components & Section-role classification; evidence retrieval &
Section-role macro-F1 0.9995; retrieval MRR 0.962, P@1 0.933 \\
\end{longtable}
}

\textbf{Table 2. Appraisal six-domain component across generations.} All
labels are deterministic weak labels unless stated; metrics measure
consistency with the stated reference, not external validity. Sources:
\texttt{docs/architecture/appraisal-step-component-results-2026-08-18.md},
\texttt{appraisal-rubric-v2-results-2026-08-18.md},
\texttt{val2c-scoring-result-2026-08-18.md}.

{\def\LTcaptype{none} % do not increment counter
\begin{longtable}[]{@{}
  >{\raggedright\arraybackslash}p{(\linewidth - 6\tabcolsep) * \real{0.2500}}
  >{\raggedright\arraybackslash}p{(\linewidth - 6\tabcolsep) * \real{0.2500}}
  >{\raggedright\arraybackslash}p{(\linewidth - 6\tabcolsep) * \real{0.2500}}
  >{\raggedright\arraybackslash}p{(\linewidth - 6\tabcolsep) * \real{0.2500}}@{}}
\toprule\noalign{}
\begin{minipage}[b]{\linewidth}\raggedright
Generation / evaluation
\end{minipage} & \begin{minipage}[b]{\linewidth}\raggedright
Reference
\end{minipage} & \begin{minipage}[b]{\linewidth}\raggedright
Metric
\end{minipage} & \begin{minipage}[b]{\linewidth}\raggedright
Value
\end{minipage} \\
\midrule\noalign{}
\endhead
\bottomrule\noalign{}
\endlastfoot
V3 (rule weak labels) & Rule labels on dev & macro-F1 & 0.8500 \\
Rubric V2 & Rubric-grounded labels on dev & macro-F1 / weighted-F1 &
0.3777 / 0.871 \\
Generation-2 spot check (n=100) & Rubric-grounded labels & agreement /
weighted-F1 & 0.80 / 0.871 \\
VAL-2c (n=100 criteria) & Guideline criteria & Cohen's kappa (95\% CI) &
0.311 (0.191 to 0.431) \\
Criterion-supervised relabeling & Guideline criteria & Records
(other/selection/reporting/performance/attrition/detection) & 9,906
(8,443/617/488/145/112/101) \\
\end{longtable}
}

\textbf{Table 3. Reconstruction of Hodgkiss et al.~2023 {[}18{]}:
deterministic R pipeline versus published values.} Three locked runs
produced byte-identical outputs (evidence:
\texttt{validation-output/reconstruction-runs/sci-exercise/rep-1/execution-receipt.json}).

{\def\LTcaptype{none} % do not increment counter
\begin{longtable}[]{@{}
  >{\raggedright\arraybackslash}p{(\linewidth - 8\tabcolsep) * \real{0.2000}}
  >{\raggedright\arraybackslash}p{(\linewidth - 8\tabcolsep) * \real{0.2000}}
  >{\raggedright\arraybackslash}p{(\linewidth - 8\tabcolsep) * \real{0.2000}}
  >{\raggedright\arraybackslash}p{(\linewidth - 8\tabcolsep) * \real{0.2000}}
  >{\raggedright\arraybackslash}p{(\linewidth - 8\tabcolsep) * \real{0.2000}}@{}}
\toprule\noalign{}
\begin{minipage}[b]{\linewidth}\raggedright
Dimension
\end{minipage} & \begin{minipage}[b]{\linewidth}\raggedright
Reconstructed
\end{minipage} & \begin{minipage}[b]{\linewidth}\raggedright
Published
\end{minipage} & \begin{minipage}[b]{\linewidth}\raggedright
Tolerance
\end{minipage} & \begin{minipage}[b]{\linewidth}\raggedright
Verdict
\end{minipage} \\
\midrule\noalign{}
\endhead
\bottomrule\noalign{}
\endlastfoot
Pooled MD (mL/kg/min) & 2.865 & 2.9 & ±0.05 & Pass (Δ −0.035) \\
CI lower & 1.795 & 1.8 & ±0.05 & Pass (Δ −0.005) \\
CI upper & 3.935 & 3.9 & ±0.05 & Pass \\
I² & 92.67\% & 93\% & ±1.0 pp & Pass (Δ −0.33) \\
Studies k & 16 & 16 & Exact & Pass \\
Egger test p & 0.483 & 0.54 & Significance side & Pass (both
non-significant) \\
\end{longtable}
}

\textbf{Table 4. Cross-provider comparison on the same blinded set of
999 paragraphs (evidence:
\texttt{docs/architecture/glm-cross-provider-results-2026-08-18.md}).}

{\def\LTcaptype{none} % do not increment counter
\begin{longtable}[]{@{}
  >{\raggedright\arraybackslash}p{(\linewidth - 6\tabcolsep) * \real{0.2500}}
  >{\raggedright\arraybackslash}p{(\linewidth - 6\tabcolsep) * \real{0.2500}}
  >{\raggedright\arraybackslash}p{(\linewidth - 6\tabcolsep) * \real{0.2500}}
  >{\raggedright\arraybackslash}p{(\linewidth - 6\tabcolsep) * \real{0.2500}}@{}}
\toprule\noalign{}
\begin{minipage}[b]{\linewidth}\raggedright
Metric
\end{minipage} & \begin{minipage}[b]{\linewidth}\raggedright
DeepSeek C3-R2
\end{minipage} & \begin{minipage}[b]{\linewidth}\raggedright
GLM glm-5.2 C3
\end{minipage} & \begin{minipage}[b]{\linewidth}\raggedright
Agreement
\end{minipage} \\
\midrule\noalign{}
\endhead
\bottomrule\noalign{}
\endlastfoot
Section-role macro-F1 & 0.9385 & 0.8816 & --- \\
Retrieval selection accuracy & 0.93 & 0.96 & --- \\
Cohen's kappa on shared 999 scored passages & --- & --- & 0.8472 (95\%
CI 0.8221 to 0.8722) \\
\end{longtable}
}

\begin{center}\rule{0.5\linewidth}{0.5pt}\end{center}

\subsection{References}\label{references}

\begin{enumerate}
\def\labelenumi{\arabic{enumi}.}
\tightlist
\item
  Higgins JPT, et al., editors. Cochrane Handbook for Systematic Reviews
  of Interventions, version 6. Cochrane; 2022.
\item
  Page MJ, et al.~PRISMA 2020. BMJ. 2021;372:n71.
\item
  Rethlefsen ML, et al.~PRISMA-S. Syst Rev.~2021;10:39.
\item
  Elliott JH, et al.~Living systematic reviews. J Clin Epidemiol. 2017.
\item
  Lu C, et al.~The AI Scientist. arXiv:2408.06292. 2024.
  doi:10.48550/arXiv.2408.06292.
\item
  Yamada Y, et al.~AI Scientist-v2. arXiv:2504.08066. 2025.
  doi:10.48550/arXiv.2504.08066.
\item
  Gottweis A, et al.~Towards an AI co-scientist. arXiv:2502.18864. 2025.
  doi:10.48550/arXiv.2502.18864.
\item
  Schmidgall S, et al.~Agent Laboratory. arXiv:2501.04227. 2025.
  doi:10.48550/arXiv.2501.04227.
\item
  Buehler MJ. PRefLexOR. npj. 2025. doi:10.1038/s44387-025-00003-z.
\item
  Shinn N, et al.~Reflexion. arXiv:2303.11366. 2023.
  doi:10.48550/arXiv.2303.11366.
\item
  Setlur A, et al.~Rewarding Progress: Scaling Automated Process
  Verifiers. In: International Conference on Learning Representations
  (ICLR); 2025.
\item
  Holub M, et al.~Reflecting in the reflection: integrating a Socratic
  questioning framework into automated AI-based question generation.
  arXiv. 2026.
\item
  Xie A, Su W, Zhou Y, Liu Y, Zhang M, Ai Q. MetaSyn: a benchmark for
  LLM agents on meta-analysis articles from Nature Portfolio.
  arXiv:2606.17041v6. 2026. doi:10.48550/arXiv.2606.17041.
\item
  Wang Y. FirstResearch: auditable question formation for LLM scientific
  discovery agents. arXiv:2607.05682. 2026.
  doi:10.48550/arXiv.2607.05682.
\item
  Zheng L, et al.~Judging LLM-as-a-judge with MT-Bench and chatbot
  arena. arXiv:2306.05685. 2023. doi:10.48550/arXiv.2306.05685.
\item
  Landis JR, Koch GG. The measurement of observer agreement for
  categorical data. Biometrics. 1977;33(1):159-174. doi:10.2307/2529310.
\item
  Cohen J. A coefficient of agreement for nominal scales. Educ Psychol
  Meas. 1960;20(1):37-46. doi:10.1177/001316446002000104.
\item
  Hodgkiss DD, Bhangu GS, Lunny C, Jutzeler CR, Chiou S-Y, Walter M, et
  al.~Exercise and aerobic capacity in individuals with spinal cord
  injury: a systematic review with meta-analysis and meta-regression.
  PLoS Med. 2023;20(11):e1004082. doi:10.1371/journal.pmed.1004082.
\item
  Brümmer LE, Katzenschlager S, Gaeddert M, Erdmann C, Schmitz S, Bota
  M, et al.~Accuracy of novel antigen rapid diagnostics for SARS-CoV-2:
  a living systematic review and meta-analysis. PLoS Med.
  2021;18(8):e1003735. doi:10.1371/journal.pmed.1003735.
\item
  Brümmer LE, Katzenschlager S, McGrath, et al.~Accuracy of rapid
  point-of-care antigen-based diagnostics for SARS-CoV-2: an updated
  systematic review and meta-analysis with meta-regression analyzing
  influencing factors. PLoS Med. 2022;19(5):e1004011.
  doi:10.1371/journal.pmed.1004011.
\item
  The PLOS Medicine Staff. Correction: accuracy of novel antigen rapid
  diagnostics for SARS-CoV-2: a living systematic review and
  meta-analysis. PLoS Med. 2021;18(10):e1003825.
  doi:10.1371/journal.pmed.1003825.
\end{enumerate}

\begin{center}\rule{0.5\linewidth}{0.5pt}\end{center}

\emph{Manuscript draft v1.0. Prepared with AI assistance under the
project's AI-use policy; all figures trace to execution receipts listed
in Data Availability. Citation order follows first appearance
(Vancouver).}

\end{document}
